\chapter{Options Markets in Europe and Asia}\label{ch:m1:options-markets-europe-asia}

In 2024 more than four of every five futures and options contracts traded on
the world's exchanges changed hands in Asia-Pacific; a year later, under
two in three, and the world total had fallen by four tenths while futures
volume rose. No crisis had occurred. A regulator had decided that a contract
which individuals were buying in enormous numbers was too small and expired
too often, and rewrote its specification. Twelve years earlier another
regulator had done the same thing to what was then the most traded option in
the world, for the same reason. Options markets outside the
United States are national, each with one dominant exchange, its own
contract design, its own retail population and its own regulator's view of
both. This chapter gives the tools for comparing them, which begin with
refusing to count contracts.

\section{Three ways to measure a market}

\begin{definition}[Lot size, notional and premium turnover]\label{def:m1:options-markets-europe-asia:turnover}
The \emph{lot size}\index{lot size (derivatives)} of a derivative is the
number of units of underlying per contract: another name for the multiplier,
used where contracts are specified in shares or index units per lot.
\emph{Notional turnover}\index{notional turnover} is contracts times lot size
times the underlying's price. \emph{Premium turnover}\index{premium turnover}
is, for options, contracts times lot size times the option's price: the money
that actually changed hands.
\end{definition}

The three measures answer different questions. Contracts measure messages
and exchange fees. Notional measures the exposure that changed hands, and
overstates it for options, most of which are far out of the money. Premium
measures the money at stake and the revenue pool available to \omterm{def:m1:what-a-trading-firm-does:market-maker}{market makers},
whose spread is a fraction of the premium.
\Cref{fig:m1:options-markets-europe-asia:rankings} ranks four invented markets,
in realistic proportions, three ways.

\begin{omfigure}
\begin{tikzpicture}
\begin{axis}[width=11cm,height=5.2cm,ybar,bar width=10pt,
  ylabel={share (\%)},
  symbolic x coords={M1,M2,M3,M4},xtick=data,
  xticklabels={{M1\\small weekly},{M2\\large index},{M3\\mid-size index},{M4\\index futures}},
  tick label style={font=\small},label style={font=\small},xticklabel style={font=\footnotesize,align=center},
  ymin=0,ymax=112,ymajorgrids,grid style={gray!25},enlarge x limits=0.15,
  nodes near coords,nodes near coords style={font=\scriptsize,/pgf/number format/.cd,fixed,precision=0},
  legend columns=3,legend style={font=\footnotesize,at={(0.5,-0.3)},anchor=north,draw=none,fill=none,/tikz/every even column/.append style={column sep=8pt}},
  table/col sep=comma]
\addplot[fill=omExo!50,draw=omExo] table[x=market,y=contracts]{figdata/markets-1/27-options-markets-europe-asia/rankings.csv};
\addplot[fill=omDef!70,draw=omDef] table[x=market,y=notional]{figdata/markets-1/27-options-markets-europe-asia/rankings.csv};
\addplot[fill=omThm!60,draw=omThm] table[x=market,y=premium]{figdata/markets-1/27-options-markets-europe-asia/rankings.csv};
\legend{by contracts,by notional,by premium}
\end{axis}
\end{tikzpicture}

\omcaption{Four invented markets. M1, weekly index options with a small lot,
has 98\% of the contracts, 48\% of the notional and 38\% of the premium; M2, a
large-contract index option, has under 1\% of the contracts and 55\% of the
premium. Data: the tutorial.}\label{fig:m1:options-markets-europe-asia:rankings}
\end{omfigure}

\begin{omfigure}
\begin{tikzpicture}
\begin{axis}[width=10.5cm,height=4.8cm,xbar,bar width=6pt,
  xlabel={contracts (billion)},
  symbolic y coords={Other,Europe,Latin America,North America,Asia-Pacific},ytick=data,
  tick label style={font=\small},label style={font=\small},
  xmin=0,xmax=200,xmajorgrids,grid style={gray!25},enlarge y limits=0.14,
  nodes near coords,nodes near coords style={font=\scriptsize,fill=white,inner sep=1pt,/pgf/number format/.cd,fixed,fixed zerofill,precision=1},
  legend style={font=\footnotesize,at={(0.97,0.5)},anchor=east,draw=none,fill=white},reverse legend,
  table/col sep=comma]
\addplot[fill=omDef!60,draw=omDef] table[y=region,x=y2025]{figdata/markets-1/27-options-markets-europe-asia/regions.csv};
\addplot[fill=omExo!40,draw=omExo] table[y=region,x=y2024]{figdata/markets-1/27-options-markets-europe-asia/regions.csv};
\legend{2025,2024 (derived)}
\end{axis}
\end{tikzpicture}

\omcaption{Listed derivatives volume by region. The 2025 figures are the
industry association's; 2024 is derived from its published year-on-year
changes. Asia-Pacific went from 170 to 76 billion contracts while every other
region grew. Data: \cref{dat:m1:index-and-single-stock-futures-worldwide:fia}.}\label{fig:m1:options-markets-europe-asia:regions}
\end{omfigure}

\section{Europe}

European equity options are concentrated on two exchange groups, each the
home of its national indices and shares, each clearing in its own house
(\cref{ch:m1:option-contracts-and-exchanges}). The flagship is the option on
the euro-area blue-chip index: European-style, \omterm{def:m1:basis-roll-and-delivery:delivery}{cash-settled}, \euro10 a point
like its future, hedged with that future and quoted against it. Three
features distinguish these markets from the American one. A large share of
institutional volume is negotiated bilaterally and registered with the
exchange as a block, so that the screen shows a fraction of the activity.
Retail participation in listed options is small; individuals who want
leverage buy bank-issued warrants and certificates, a parallel market with
its own venues and its own \omterm{def:m1:what-a-trading-firm-does:market-maker}{market makers}, the issuers themselves. And there
is one order book per product, so that the routing and auction machinery of
\cref{ch:m1:options-market-structure} has no equivalent: the structure is
that of a futures exchange.

\section{Korea}

In the 2000s the most traded derivative in the world, by contracts, was the
option on the Korean blue-chip index: a small contract, cheap
out-of-the-money strikes, and a very large population of individual
traders.

\begin{dated}{2012-06}{The multiplier, times five}\label{dat:m1:options-markets-europe-asia:kospi}
On 14 June 2012 the contract size of the index option was raised from 100\,000
to 500\,000 won per point, on the home exchange and on the European exchange
that listed a linked product.
\end{dated}

The change multiplied the price of the cheapest lottery ticket by five.
Contract counts must fall fivefold if activity is unchanged; how much
activity itself falls is the question of exercise 3. It is the
template for what followed elsewhere, and a standing warning for any firm
whose business depends on one market's retail flow: the contract
specification is a policy instrument.

\section{India}

\begin{definition}[Weekly expiry and position limit]\label{def:m1:options-markets-europe-asia:weekly}
A \emph{weekly expiry}\index{weekly expiry} is an option expiry listed each
week on a fixed weekday, in addition to the monthly cycle. A \emph{position
limit}\index{position limit} is a cap set by an exchange or regulator on the
size of the position one participant or group may hold in a contract,
measured on a stated basis (gross or net, in contracts, notional or
delta-adjusted) and monitored at stated times.
\end{definition}

By 2024 the Indian index options market was the largest in the world by
contracts, by an order of magnitude: small lots, several indices each with a
\omterm{def:m1:options-markets-europe-asia:weekly}{weekly expiry} on a different weekday, so that some index expired every day,
and millions of individuals trading options on their expiry day.

\begin{dated}{2026-09}{What the regulator found, and did}\label{dat:m1:options-markets-europe-asia:sebi}
\emph{The study.} The securities regulator's study published on 23 September
2024 found that 93\% of individual traders in equity derivatives lost money
over the three financial years to March 2024, with aggregate losses above
1.8 trillion rupees.
\emph{The measures.} Its circular of 1 October 2024 set the contract value of
index derivatives at 1.5 to 2 million rupees when reviewed (from 20 November
2024), allowed each exchange weekly expiries on only one benchmark index,
required \omterm{def:m1:option-contracts-and-exchanges:option}{option premiums} to be collected upfront from buyers, removed the
margin benefit of \omterm{def:m1:matching-algorithms-and-implied-spreads:calendar}{calendar spreads} on expiry day, added a 2\% extreme-loss
margin on short options on expiry day, and introduced intraday monitoring of
\omterm{def:m1:options-markets-europe-asia:weekly}{position limits}.
\emph{The order.} On 3 July 2025 the regulator issued an interim order against
a foreign proprietary trading group, alleging manipulation of a bank index on
expiry days between January 2023 and 2025, and impounded 48.4 billion rupees
of alleged gains. The group deposited the amount on 14 July 2025 while
reserving its right to contest the order. The matter was not decided at the
time of writing.
\end{dated}

\section{The expiry-day case}

Set aside the question of what one firm did, which a tribunal will answer,
and consider the structure that the order describes, because it is general.
An index has a few dozen constituents whose cash and futures markets have a
certain depth. Options on the index, expiring today, have a \omterm{def:m1:options-markets-europe-asia:turnover}{notional turnover}
many times that depth. The options settle on the index's closing value,
computed from the constituents' prices over the final minutes. Then a
participant large enough to move the constituents by a fraction of a percent
changes the payoff of an options position many times larger than the shares
traded.

\begin{proposition}[Leverage of the settlement]\label{prop:m1:options-markets-europe-asia:leverage}
Let a participant hold index options with aggregate delta-one-equivalent
exposure $N$ at expiry (the notional of the options that finish in the
money), and let moving the index's settlement value by a fraction
$\varepsilon$ cost $c(\varepsilon)$ in impact on the underlying shares. The
payoff changes by $N\varepsilon$. With the square-root rule,
$c(\varepsilon) \propto \varepsilon^{3}$ for the shares traded to produce
$\varepsilon$, so for small $\varepsilon$ the gain is linear and the cost
cubic: whenever $N$ is large relative to the depth of the underlying, some
$\varepsilon > 0$ is profitable.
\end{proposition}
\begin{proof}
Under \cref{met:m1:the-sell-side:sqrt}, impact $\varepsilon \propto
\sigma\sqrt{Q/V}$, so the quantity needed is $Q \propto V\varepsilon^2/\sigma^2$
and the cost of trading it, quantity times average impact, is proportional to
$Q\varepsilon \propto \varepsilon^3$.
\end{proof}

This is why settlement procedures use averages over a window, why regulators
monitor expiry-day positions intraday, and why a firm's compliance function
must look at its cash and derivatives books \emph{together}: trades that are
each defensible as hedging or arbitrage can, in combination and with intent,
constitute manipulation, and the pattern is visible afterwards in exactly the
data the firm itself keeps. The distinction between moving a price as a
by-product of legitimate trading and trading in order to move it is drawn by
law, not by arithmetic, and differs by jurisdiction; One Quant Book 16
returns to it.

\begin{omfigure}
\begin{tikzpicture}[font=\small,
  bk/.style={draw=omDef,rounded corners=2pt,minimum width=3.4cm,minimum height=1.1cm,align=center,fill=omDef!4}]
\node[bk] (und) at (0,0) {Constituent shares\\and their futures};
\node[bk,draw=omThm,fill=omThm!5] (opt) at (8.2,0) {Index options\\expiring today};
\node[bk,draw=omExo,fill=omExo!6,minimum width=3.0cm] (idx) at (4.1,2.1) {Index settlement value};
\draw[-{Stealth[length=2mm]},thick,omDef] (und.north) |- node[pos=0.25,left,font=\footnotesize,align=right]{prices over the\\final window} (idx.west);
\draw[-{Stealth[length=2mm]},thick,omThm] (idx.east) -| node[pos=0.75,right,font=\footnotesize,align=left]{payoff of every\\in-the-money option} (opt.north);
\node[font=\footnotesize,below=4pt of und] {depth: $V$};
\node[font=\footnotesize,below=4pt of opt] {exposure: $N \gg V$};
\end{tikzpicture}

\omcaption{The structure behind every settlement-manipulation case: a small
market whose prices determine the payoff of a much larger
one.}\label{fig:m1:options-markets-europe-asia:structure}
\end{omfigure}

\section{Other Asian markets}

Elsewhere in Asia the picture is closer to Europe's: one derivatives exchange
per country, index options as the main product, \omterm{def:m1:option-contracts-and-exchanges:style}{European exercise} and \omterm{def:m1:basis-roll-and-delivery:delivery}{cash
settlement}, futures-style market structure. Japan's and Hong Kong's index
options are dominated by institutions and by the hedging of structured
products sold to individuals, whose issuers are structurally long or short
volatility in ways that shape the local surface; Taiwan's index options have
a large retail share. For a trading firm the questions of
\cref{ch:m1:asian-and-emerging-equity-markets} come first in each: who may
trade, through whom, with what \omterm{def:m1:options-markets-europe-asia:weekly}{position limits}, under what tax.

\section{Tutorial: three rankings}\label{tut:m1:options-markets-europe-asia:rankings}

\begin{tutorial}
\textbf{Goal.} Convert contract counts into notional and \omterm{def:m1:options-markets-europe-asia:turnover}{premium turnover} in
one currency, rank markets three ways, and compute what a lot-size change
should do to a contract count. \textbf{End state:} the two data figures of
this chapter.
\begin{tutsteps}
\item \textbf{One record per market.}
\omcode{code/firm/turnover/firm_turnover.py}{9}{30}{Contracts, lot, underlying, premium, currency: three measures of turnover.}\label{lst:m1:options-markets-europe-asia:stats}
\item \textbf{Lot changes.} If notional activity is unchanged, tripling the
lot divides the count by three; the count observed afterwards tells how much
activity was lost.
\omcode{code/firm/turnover/firm_turnover.py}{46}{54}{The arithmetic of a lot-size change, in both directions.}\label{lst:m1:options-markets-europe-asia:lot}
\end{tutsteps}
\textbf{What to change next.} Replace the average premium by a distribution
over strikes and estimate a \omterm{def:m1:what-a-trading-firm-does:market-maker}{market maker}'s revenue pool as half the quoted
spread times the \omterm{def:m1:options-markets-europe-asia:turnover}{premium turnover}, by market.
\end{tutorial}

\section{Build: the turnover normaliser}\label{bld:m1:options-markets-europe-asia:turnover}

\begin{build}
\bfield{Purpose} When the miniature firm's strategy committee asks which
market to enter next, the first slide must not be a league table of
contracts.
\bfield{Interface} \texttt{MarketStats(name, kind, contracts, multiplier,
underlying, avg\_premium, fx\_to\_usd)} with \texttt{notional\_usd},
\texttt{premium\_usd}, \texttt{premium\_to\_notional\_bp};
\texttt{shares(markets, measure)}; \texttt{ranking(markets, measure)};
\texttt{contracts\_after\_lot\_change(\ldots)}; \texttt{activity\_kept(\ldots)}.
\bfield{Rules} One currency for every comparison, with the exchange rate
stored beside the data. Futures have no \omterm{def:m1:options-markets-europe-asia:turnover}{premium turnover} and are excluded,
with an error, from a premium ranking of futures alone. Ties are broken by
name.
\bfield{Acceptance tests} \texttt{code/firm/turnover/tests/}: units; three
measures giving three rankings; the lot-change arithmetic both ways.
\bfield{Stretch} Delta-adjusted notional for options, from a strike
distribution of volume.
\end{build}

\begin{omsources}
\item Securities and Exchange Board of India, press release of 23 September
2024 on the updated study of individual traders in equity derivatives;
circular SEBI/HO/MRD/TPD-1/P/CIR/2024/132 of 1 October 2024.
\item Securities and Exchange Board of India, interim order of 3 July 2025 in
the matter of index manipulation; commentary: ``Jane Street and the Expiry
Day Trap'', \emph{Oxford Business Law Blog}, July 2025.
\item Eurex, ``Eurex and Korea Exchange further expand their link'' (on the
June 2012 change of contract size).
\item FIA, \emph{ETD Volume --- December 2025}.
\end{omsources}

\section{Exercises}

\begin{exercise}[$\star$]\label{exo:m1:options-markets-europe-asia:1}
An index option has a lot of 25, the index is at 22\,000 rupees and the
option costs 35 rupees per unit. With the rupee at \$0.012, give the notional
and the premium of one contract in dollars.
\end{exercise}

\begin{exercise}[$\star$]\label{exo:m1:options-markets-europe-asia:2}
Same for an index option with a multiplier of 100, the index at 5\,600 dollars
and a premium of 16. How many contracts of the previous exercise make the
premium of one of these?
\end{exercise}

\begin{exercise}[$\star$]\label{exo:m1:options-markets-europe-asia:3}
A contract trades 90 billion times a year with a lot of 25. The lot is raised
to 75. What count should be expected if notional activity is unchanged? The
count turns out to be 24 billion. What happened to activity?
\end{exercise}

\begin{exercise}[$\star\star$]\label{exo:m1:options-markets-europe-asia:4}
From \cref{fig:m1:options-markets-europe-asia:regions} compute Asia-Pacific's
share of world contracts in 2024 and in 2025, and the world total in 2024.
Check the total against the published fall of 42.2\%.
\end{exercise}

\begin{exercise}[$\star\star$]\label{exo:m1:options-markets-europe-asia:5}
A \omterm{def:m1:what-a-trading-firm-does:market-maker}{market maker} earns on average a quarter of the quoted spread, and spreads
are 4\% of premium in market M1 and 1.5\% in M2. Using \omterm{def:m1:options-markets-europe-asia:turnover}{premium turnover} of
\$945 billion and \$1\,360 billion, estimate the revenue pool in each. Which
market would you enter, and what else would you need to know?
\end{exercise}

\begin{exercise}[$\star\star$]\label{exo:m1:options-markets-europe-asia:6}
Give two reasons why a settlement value computed as an average over the last
thirty minutes is harder to move than a last-trade price, and one cost of
using it.
\end{exercise}

\begin{exercise}[$\star\star\star$]\label{exo:m1:options-markets-europe-asia:7}
\emph{Coding.} With the four markets of the tutorial, print the three tables
of shares. Then change M1's lot to 75 with \texttt{activity\_kept} of 0.8 and
print them again. Which shares move?
\end{exercise}

\begin{exercise}[$\star\star\star$]\label{exo:m1:options-markets-europe-asia:8}
\emph{Find the flaw.} A desk head: ``In the morning we bought bank stocks
because they were cheap to futures. In the afternoon we sold them because our
risk limit required it. Separately, our options desk was short the index.
Each trade has its own rationale, so there is nothing to explain.'' What is
missing from this defence?
\end{exercise}

\section{Problem: Counting Contracts}

\begin{problem}[{Weekend problem --- where is the money?}]\label{pb:m1:options-markets-europe-asia:1}
A firm compares two index option markets for a year. M1: 90 billion
contracts, lot 25, index 22\,000 rupees, average premium 35 rupees, rupee at
\$0.012, spreads 4\% of premium, exchange and regulatory fees 0.05\% of \omterm{def:m1:options-markets-europe-asia:turnover}{premium
turnover}, a transaction tax on option sales of 0.1\% of premium. M2: 850
million contracts, multiplier 100, index \$5\,600, average premium \$16,
spreads 1.5\% of premium, fees \$0.50 a contract for a \omterm{def:m1:what-a-trading-firm-does:market-maker}{market maker}.

\medskip\noindent\textbf{Part I --- Three measures.}
\begin{enumerate}
\item Give the ratio of contracts, M1 to M2.
\item Give each market's \omterm{def:m1:options-markets-europe-asia:turnover}{notional turnover} in dollars.
\item Give each market's \omterm{def:m1:options-markets-europe-asia:turnover}{premium turnover} in dollars.
\item Give the \omterm{def:m1:options-markets-europe-asia:turnover}{premium turnover} ratio, M2 to M1.
\item Give premium as basis points of notional in each. What does the
difference say about what is being traded?
\end{enumerate}

\medskip\noindent\textbf{Part II --- A revenue pool.}
\begin{enumerate}[resume]
\item If \omterm{def:m1:what-a-trading-firm-does:market-maker}{market makers} collectively earn a quarter of the spread on all
premium, give each market's gross pool.
\item Give M1's tax and fees on the makers' side, assuming makers are on one
side of every trade and sell half the time.
\item Give M2's fees for makers on one side of every trade.
\item Give the net pools.
\item What fraction of each pool could a new entrant realistically win, and
what determines it?
\end{enumerate}

\medskip\noindent\textbf{Part III --- Policy risk.}
\begin{enumerate}[resume]
\item M1's regulator triples the lot and removes four of five weekly
expiries. Contracts fall to 24 billion. Give the notional activity kept.
\item Give the new \omterm{def:m1:options-markets-europe-asia:turnover}{premium turnover} if the average premium per unit is
unchanged.
\item Recompute M1's net pool.
\item The firm had hired twelve people and rented co-location for M1. What
is the lesson for the business plan?
\end{enumerate}

\medskip\noindent\textbf{Part IV --- Judgement.}
\begin{enumerate}[resume]
\item Why did M1's regulator act? Use
\cref{dat:m1:options-markets-europe-asia:sebi}.
\item Who was on the other side of the individuals' losses?
\item Does that make market making in M1 illegitimate? Where is the line
that \cref{prop:m1:options-markets-europe-asia:leverage} points to?
\item A foreign firm entering M1 needs what, besides a strategy?
\item State the \emph{named result}: the \omterm{def:m1:options-markets-europe-asia:turnover}{premium turnover} ratio, M2 to M1.
\item In one sentence: what does a league table of contracts measure?
\end{enumerate}
\end{problem}

\section{Interview questions}

\begin{interviewq}[$\star$ \iqroles{trader, researcher}]\label{iq:m1:options-markets-europe-asia:1}
An exchange claims to be ``the world's largest derivatives exchange''. What
do you ask?
\end{interviewq}

\begin{interviewq}[$\star$ \iqroles{trader, researcher, bank}]\label{iq:m1:options-markets-europe-asia:2}
How do listed options markets in Europe differ from those in the United
States?
\end{interviewq}

\begin{interviewq}[$\star\star$ \iqroles{researcher, trader}]\label{iq:m1:options-markets-europe-asia:3}
A regulator triples the minimum contract size of index options. Predict the
effects on volume, on participants and on \omterm{def:m1:what-a-trading-firm-does:market-maker}{market makers}' revenue.
\end{interviewq}

\begin{interviewq}[$\star\star$ \iqroles{trader, researcher}]\label{iq:m1:options-markets-europe-asia:4}
Why are \omterm{def:m1:basis-roll-and-delivery:delivery}{cash-settled} index derivatives vulnerable around their settlement,
and how do exchanges defend against it?
\end{interviewq}

\begin{interviewq}[$\star\star$ \iqroles{trader, bank}]\label{iq:m1:options-markets-europe-asia:5}
Your cash desk and your options desk are both active in the same index on
expiry day. What controls do you want?
\end{interviewq}

\begin{interviewq}[$\star\star\star$ \iqroles{researcher, trader}]\label{iq:m1:options-markets-europe-asia:6}
You are asked to evaluate entering a fast-growing retail options market
abroad. Structure the analysis.
\end{interviewq}
